%
% 11-problem.tex -- Problemstellung
%
% (c) 2026 Prof Dr Andreas Müller
%
\subsection{Problemstellung}
Die von Bernoulli gelöste Aufgabe ist die Stammfunktion einer
rationalen Funktion zu finden, für die ausserdem angenommen wird,
dass eine Faktorisierung des Nenners in Linearfaktoren bekannt ist.
Der Fundamentalsatz der Algebra garantiert, dass eine solche
Faktorisierung immer exisiert, wenngleich er keine Methode
bereitstellt, die Faktorisierung zu finden.

\begin{aufgabe}
Gegeben ist eine rationale Funktion $f\in \mathbb{C}(z)$ mit
komplexen Koeffizienten.
Finde eine Stammfunktion $F$.
\end{aufgabe}

Aufgrund seiner Annahmen kann Bernoulli annehmen, dass zur
gegebenen rationale Funktion eine Partialbruchzerlegung 
wie in Abschnitt~\ref{buch:analytisch:rational:subsection:partialbruch}
gefunden werden kann.
Der Integrand kann daher als Summe
\begin{align*}
f(z)
=
p(z)
+
\sum_{k=1}^n\sum_{l=1}^{n_k}\frac{A_k^l}{(z-\alpha_k)^l}
\end{align*}
geschrieben werden, wobei $p(z)$ ein Polynom und $\alpha_k$ die Nullstellen
des Nenners von $f(z)$ sind. 
Um eine Stammfunktion zu finden, müssen Stammfunktionen der
einzelnen Summanden gefunden werden.
Der Polynomterm ist dabei das kleinste Problem, da die einzelnen
Summanden $z^n$ die Stammfunktion $z^{n+1}/(n+1)$ haben.
Etwas mehr Arbeit geben die Brüche.
